% Section "learn" (writer: learn). Sources: research/tracks/foundations/notes-final.md §6 (Thms 6.1, 6.5, 6.7,
% Cor 6.6, Props 6.8--6.10, Rem 6.11, §6.6, §6.7) and research/tracks/context/notes-final.md §2.4, §5.2, §5.4.
% Proofs and the secondary items prop:learn:active, rem:learn:increasing, rem:learn:cesaro, rem:learn:relations,
% rem:learn:positive, rem:learn:interventions are in sections/app-learn.tex, to respect the 5-page cap (OUTLINE.md).
\section{Learning theorems: SFI as supervised learning}\label{sec:learn}

H\"anni describes SFI as ``just ideal supervised learning'' \citep{hanni2026request}. This section states the guarantees behind that description and where they stop. Throughout, $D_n=\{(x_t,y_t):t\le n\}$ and $\LSFI(D_n)=-\log\MF(D_n)$ is SFI's cumulative log loss in the semimeasure convention, independent of order and repetitions and infinite on inconsistent sequences (\cref{prop:setting:order}); upper bounds in it also hold normalised (\cref{prop:setting:normalise}). We write $\mu$ for the input distribution, $c$ for a machine constant, and $L_\nu(D):=-\log\nu(D)$. For an lsc conditional semimeasure $P$ (\cref{prop:setting:includes}(a)), $K(P)$ is the least $K(i)$ over its indices in the enumeration of \cref{def:versions:variants}. For a partial computable $\varphi$, $K(\varphi)$ is the least $K(e)$ with $\varphi_e=\varphi$.

% ------------------------------------------------------------------
\subsection{Dominance}\label{sec:learn:dominance}

\begin{theorem}[dominance]\label{thm:learn:dominance}
\status{proved}\src{foundations Thm 6.1}
Let $\nu$ be a machine law (\cref{def:setting:machine}). Then $\MF(D)\ge2^{-K(\nu)}\nu(D)$ for every data set $D$. Hence for every consistent sequence, however its inputs are chosen (in advance, adversarially, adaptively, non-computably),
\begin{equation}\label{eq:learn:dominance}
\LSFI(D_n)\ \le\ -\log\nu(D_n)+K(\nu).
\end{equation}
In particular:
(a) \emph{x-only samplers:} for every lsc conditional semimeasure $P$, $\LSFI(D_n)\le\sum_{t\le n}-\log P(y_t\mid x_t)+K(P)+c$, where the right side counts a repeated input each time it appears and the left side once;
(b) \emph{deterministic programs:} if $y_t=\varphi(x_t)$ for all $t$, with $\varphi$ partial computable, then $\LSFI(D_n)\le K(\varphi)+c$;
(c) \emph{weak comparators:} for every lsc weak semimeasure $\nu$ (\cref{def:setting:weak}), $\LSFI(D_n)\le-\log\nu(D_n)+K(\nu)+K(\dom D_n)+c$. The term $K(\dom D_n)$ cannot be replaced by a constant (semimeasure convention): on the $2k$-point data sets of \cref{thm:setting:weakvsmachine}(iii), $\LSFI$ exceeds $-\log\Mw$ by at least $k\log\frac32-K(k)-c$, while $K(\dom D)\le k\log3+K(k)+c$.
\end{theorem}

\begin{proof}
\Cref{thm:setting:universal} gives $\MF\ge2^{-K(\nu)}\nu$, and \cref{prop:setting:order} gives $\LSFI(D_n)=-\log\MF(D_n)$ in every order. (a) The product law $\nu_P(D)=\prod_{(x,y)\in D}P(y\mid x)$ is a machine law with $K(\nu_P)\le K(P)+c$ (\cref{prop:setting:includes}(a)); it counts each input once, and the extra factors of repeats are at most~$1$. (b) \Cref{prop:setting:includes}(b). (c) \Cref{thm:setting:weakvsmachine}(ii),(iii); each domain of that family is determined by $k$ and a vector in $\{1,2,3\}^k$.
\end{proof}

Taking expectations gives the bound under every data distribution. Part (a) is H\"anni's ``Thm 1 (template)'' for computably samplable x-only predictors \citep{hanni2026notes}: each ``loses to `the solomonoff inductor for C'{}'' by ``at most a constant (depending on $P$ but not the distribution)''. It holds in his setting of ``input-output pairs from a function'' (consistent sequences), with randomness read once per input, as his remark on ``different bits of randomness \dots{} for sampling on different inputs'' suggests.

\begin{remark}[scope of dominance]\label{rem:learn:h2}
\status{proved; open part}\src{foundations §6.1, F20--F23}
(i) The comparators include all computable laws of random total functions, for every countable $Y$ (\cref{prop:setting:includes,lem:setting:total}).
(ii) The label sequence must be consistent: if an input reappears with another label, $\LSFI=\infty$, while a sampler with fresh randomness per \emph{query} has finite loss. Randomness per input is the right reading for data from a function (\cref{rem:setting:perinput}); per query it gives $\Mpq$ (\cref{prop:learn:perquery}).
(iii) For general lsc laws of random partial functions the statement is open (\cref{q:setting:rpf}).
\end{remark}

% ------------------------------------------------------------------
\subsection{Merging and mistakes}\label{sec:learn:merging}

\begin{theorem}[merging]\label{thm:learn:merging}
\status{proved; finite analogue computed (foundations c4, r6; integrability inequality c11(c))}\src{foundations Thm 6.5 and remarks}
Let $\nu$ be a machine law of a random total function, equivalently a computable law of a random total function (\cref{lem:setting:total}), and $F\sim\nu$. Let the inputs follow any measurable rule $x_t=s_t(h_{t-1},\xi_t)$, with $h_{t-1}$ the history, $(\xi_t)$ independent of $F$, and $s_t$ arbitrary (possibly non-computable); $y_t:=F(x_t)$. With the one-step conditionals of \cref{def:setting:data} let
\[ d_t:=\KL\big(\nu(\cdot\mid D_{t-1},x_t)\,\big\|\,\MF(\cdot\mid D_{t-1},x_t)\big)\ \text{bits}\qquad(d_t=0\text{ at a repeated input}). \]
Then $\sum_{t\le n}\Ex d_t=\Ex\log(\nu(D_n)/\MF(D_n))\le K(\nu)$ for every $n$. Hence $\sum_{t\ge1}\Ex d_t\le K(\nu)$, $d_t\to0$ almost surely, and $\sum_{t\ge1}\Ex\|\nu(\cdot\mid D_{t-1},x_t)-\MF(\cdot\mid D_{t-1},x_t)\|_1^2\le2\ln2\cdot K(\nu)$.
\end{theorem}

\emph{Proof idea.} Given the past, $y_t$ has law $\nu(\cdot\mid D_{t-1},x_t)$, so the expected log ratio at step $t$ is $d_t$; the ratios telescope to $\log(\nu(D_n)/\MF(D_n))\le K(\nu)$ (\cref{thm:learn:dominance}; integrability in \cref{app:learn:merging}).

This is the function analogue of Solomonoff's convergence theorem \citep{solomonoff1978complexity} (source not checked). If the rule may see $F$, e.g.\ picks inputs by their label, the expected bound can fail (summed KL $12.5$ bits against a bound of $1$ in a finite analogue, foundations r6); \eqref{eq:learn:dominance} still holds, and \cref{thm:hist:directed} bounds the excess by directed information.

\begin{corollary}[mistake bound]\label{cor:learn:mistakes}
\status{proved; known (Littlestone 1988; Barzdin--Freivalds 1972, content unverified)}\src{foundations Cor 6.6}
Let $\varphi$ be partial computable, $x_1,x_2,\ldots\in\dom\varphi$ any inputs, and $y_t=\varphi(x_t)$. At most $K(\varphi)+c$ steps have $\MF(\varphi(x_t)\mid D_{t-1},x_t)\le\frac12$. So $\hat y_t:=\arg\max_y\MF(y\mid D_{t-1},x_t)$ makes at most $K(\varphi)+c$ mistakes on every input sequence.
\end{corollary}

\begin{proof}
Each such step adds at least one bit to $\LSFI(D_n)\le K(\varphi)+c$ (\cref{thm:learn:dominance}(b)); above $\frac12$, $\varphi(x_t)$ is the unique maximiser.
\end{proof}

This is the halving bound \citep{littlestone1988learning,barzdin1972prediction} with $K(\varphi)$ for $\log|H|$; the predictor is not computable (\cref{prop:setting:computability}); examples: \cref{prop:ex:deterministic}.

% ------------------------------------------------------------------
\subsection{Online-to-batch and noise}\label{sec:learn:batch}

\begin{theorem}[online-to-batch]\label{thm:learn:batch}
\status{proved; finite analogues computed (foundations r9; c7, r7)}\src{foundations Thm 6.7}
Let $x_1,x_2,\ldots$ be i.i.d.\ from any distribution $\mu$ on $X$, computable or not.
\begin{enumerate}[label=(\alph*),leftmargin=*]
\item \emph{Deterministic labels.} Let $\varphi$ be partial computable, $\mu(\dom\varphi)=1$, $y_t=\varphi(x_t)$. Put $Q_t(\cdot\mid x):=\MF(\cdot\mid D_t,x)$, $\bar Q_n:=\frac1n\sum_{t<n}Q_t$, $R(Q):=\Ex_{x\sim\mu}[-\log Q(\varphi(x)\mid x)]$. Then $\sum_{t<n}\Ex R(Q_t)\le K(\varphi)+c$, $\Ex R(\bar Q_n)\le(K(\varphi)+c)/n$, and the expected 0--1 risk $\Ex\,\mu\{x:\arg\max_y\bar Q_n(y\mid x)\ne\varphi(x)\}$ is at most $(K(\varphi)+c)/n$.
\item \emph{Noisy labels, stamped inputs.} Let $y_t\sim P(\cdot\mid x_t)$ independently given the inputs, for any conditional probability $P$ (possibly non-computable). Run SFI on $\tilde x_t:=\langle x_t,t\rangle$, for a computable bijection $\langle\cdot,\cdot\rangle:X\times\mathbb N\to X$, with data sets $\tilde D_t$; put $q_t(\cdot\mid x):=\MF(\cdot\mid\tilde D_{t-1},\langle x,t\rangle)$ and $\bar q_n:=\frac1n\sum_{t\le n}q_t$. Then for every lsc conditional semimeasure $P'$,
\[ \Ex\,\Ex_{x\sim\mu}\KL\big(P(\cdot\mid x)\,\|\,\bar q_n(\cdot\mid x)\big)\le\tfrac1n\big(K(P')+c\big)+\Ex_{x\sim\mu}\KL\big(P(\cdot\mid x)\,\|\,P'(\cdot\mid x)\big). \]
For lsc (hence computable) $P$ this is $(K(P)+c)/n$. The expected excess 0--1 risk of $\arg\max_y\bar q_n(y\mid x)$ over the Bayes classifier is at most $\sqrt{2\ln2}$ times the square root of the right side.
\item \emph{Frozen noise.} If $F$ has the product law of a computable conditional probability $P$ (one draw per input) and $y_t=F(x_t)$, then $\sum_{t\le n}\Ex d_t\le K(P)+c$, with $d_t$ as in \cref{thm:learn:merging}, for every input rule that sees only the past.
\end{enumerate}
\end{theorem}

Part (a) is online-to-batch conversion \citep{littlestone1989online,cesabianchi2004generalization}: as $x_{t+1}$ is independent of $D_t$, $\Ex R(Q_t)$ is the expected online loss at step $t+1$, \cref{thm:learn:dominance}(b) bounds the sum, and convexity moves the average inside. Averaging over a random order of a finite set instead gives \cref{prop:learn:splits}; \citet{haussler1994predicting} use a permutation argument of this kind (not checked). In (b) the stamps make inputs distinct; (c) is \cref{thm:learn:merging} for a product law (proofs: \cref{app:learn:batch}). The bounds are for the Ces\`aro average $\bar Q_n$, not for $Q_n$ (\cref{sec:learn:curves}).

\begin{proposition}[random-function SFI cannot model fresh label noise]\label{prop:learn:freshnoise}
\status{proved; refuted (the brief's fresh-noise claim, as stated); computed (foundations c5, r9(b))}\src{foundations Prop 6.8; F31}
Let $\mu(x_0)=\alpha>0$, let $P(\cdot\mid x_0)$ give positive probability to two labels, and $p_{\max}:=\max_yP(y\mid x_0)<1$. Draw the pairs $(x_t,y_t)$ i.i.d.\ from $\mu\times P$, with a fresh label at every draw, and predict with $\nu(\cdot\mid D,x)$ for a weak semimeasure $\nu$, raw or normalised, with the repeated-input convention of \cref{def:setting:data}. This covers $\MF$, $W$, $\Mprod$ and $\Mw$.
(a) For every $n\ge2$, $\Prob[D_n\text{ inconsistent}]>0$, so the expected cumulative log loss is infinite, in both conventions.
(b) $\Prob[D_n\text{ consistent}]\le(1-\alpha)^n+p_{\max}^{-1}(1-\alpha+\alpha p_{\max})^n\to0$.
(c) The statement ``for every input distribution and every computable conditional probability $P$, SFI's expected cumulative log loss exceeds that of $P$ by at most $K(P)+O(1)$'' is false for random-function SFI. For $P$ with finite expected log loss it holds for all $n$ exactly when $P(\cdot\mid x)$ is a point mass for every $x$ with $\mu(x)>0$; then it is \cref{thm:learn:batch}(a).
\end{proposition}

(Proof: \cref{app:learn:batch}.) Random-function SFI models the label as a function of the input; a fresh-noise source is not one. With $10$ equally likely inputs and fair-coin labels the data stay consistent with probability $0.95$, $0.113$, $5\cdot10^{-4}$ at $n=2,10,20$ (foundations c5). This refutes the fresh-noise form of the supervised-learning hypothesis H7 for random-function SFI (\cref{tab:intro:verdicts}). The finite-loss proviso in (c) is ours (for infinite loss the inequality is trivial). The repairs are stamps and frozen noise (\cref{thm:learn:batch}(b),(c); example: \cref{prop:ex:frozen}) and mixing per query:

\begin{proposition}[the per-query mixture]\label{prop:learn:perquery}
\status{proved; computed (foundations c7, r7)}\src{foundations Prop 6.9}
Let $\Mpq(y_{1:n}\Vert x_{1:n})=\sum_i\mm(i)\prod_{t\le n}P_i(y_t\mid x_t)$ on all finite sequences of pairs (\cref{def:versions:variants}).
(a) $\Mpq$ is an lsc chronological semimeasure: $\Mpq(\epsilon\Vert\epsilon)\le1$ and $\sum_y\Mpq(y_{1:n}y\Vert x_{1:n}x)\le\Mpq(y_{1:n}\Vert x_{1:n})$. It depends only on the multiset of pairs; it equals $\Mprod(D_n)$ on sequences with distinct inputs, is at most $\Mprod(D_n)$ on consistent sequences with repeats (often strictly), and is positive on inconsistent sequences.
(b) On every sequence, consistent or not, $-\log\Mpq(y_{1:n}\Vert x_{1:n})\le\sum_{t\le n}-\log P(y_t\mid x_t)+K(P)$ for every lsc conditional semimeasure $P$.
(c) Let $(x_t,y_t)$ be i.i.d.\ from $\mu\times P$, $\mu$ arbitrary and $P$ any conditional probability, $q_t(\cdot\mid x):=\Mpq(y_{<t}\,\cdot\Vert x_{<t}x)/\Mpq(y_{<t}\Vert x_{<t})$ and $\bar q_n:=\frac1n\sum_{t\le n}q_t$. For every lsc conditional semimeasure $P'$, $\sum_{t\le n}\Ex\,\Ex_\mu\KL(P\|q_t)\le K(P')+n\,\Ex_\mu\KL(P\|P')$, so $\Ex\,\Ex_\mu\KL(P\|\bar q_n)\le K(P')/n+\Ex_\mu\KL(P\|P')$.
(d) With the stamps of \cref{thm:learn:batch}(b), $\MF(\tilde D_n)\ge2^{-c}\Mpq(y_{1:n}\Vert x_{1:n})$.
\end{proposition}

So the regret bound holds under fresh noise for $\Mpq$, with no stamps (proofs: \cref{app:learn:perquery}). As secondary sources describe it, operator induction is $\Mpq$ (\cref{rem:versions:operator}); \SFIhist{} dominates it (\cref{prop:hist:operator}); \cref{prop:comp:round} is a time-bounded analogue.

% ------------------------------------------------------------------
\subsection{Learning curves and in-distribution error}\label{sec:learn:curves}

\begin{proposition}[the expected risk of a Bayes predictor need not decrease]\label{prop:learn:curve}
\status{proved; computed (foundations c11(a), r12)}\src{foundations Prop 6.10}
Let $X=\{a,b\}$, $\mu$ uniform, target $\varphi\equiv0$, and a Bayes mixture of three total functions: $\varphi$ with prior weight $w$, and $(1,0)$ and $(0,1)$ with weight $\frac{1-w}2$ each. The expected log risk $r_t$ of the posterior predictive after $t$ i.i.d.\ samples is $r_0=-\log\frac{1+w}2$, $r_1=\frac12\log\frac{1+w}{2w}$ and $r_t=2^{1-t}r_1$ for $t\ge1$. So $r_1>r_0$ whenever $(1+w)^3>8w$, in particular for all $w\le\frac15$; for $w=\frac1{16}$, $r_0=0.912$ and $r_1=1.544$ bits. The sum $\sum_tr_t$ increases to $\log(1/w)$.
\end{proposition}

The mixture is a random total function, so both conventions agree, but it is not $\MF$: ``more data never hurts in expectation'' fails for Bayes predictors with a fixed target; for $\MF$ this is not established. $63$ of $300$ random finite instances had non-monotone curves (foundations r12); see also \cref{rem:learn:increasing}.

\begin{remark}[last iterate and leave-one-out]\label{rem:learn:lastiterate}
\status{identity proved; computed (foundations c11(b)); bound open}\src{foundations Rem 6.11, §8 item 7}
In the setting of \cref{thm:learn:batch}(a), exchangeability gives $\Ex R(Q_n)=\Ex\big[\frac1{n+1}\sum_{i\le n+1}-\log Q^{(-i)}(\varphi(x_i)\mid x_i)\big]$, where $Q^{(-i)}$ is SFI's predictive given the data set of the samples $j\ne i$ (a held-out input that occurs among the others costs~$0$). Whether $\Ex R(Q_n)\le(K(\varphi)+c)/(n+1)$ is thus a leave-one-out question (cf.\ \cref{conj:learn:loo}); it is open. In $300$ random finite instances the largest $(n+1)r_n/K$ was $0.873$ (foundations r12).
\end{remark}

\emph{Errors.} For $\nu\in\{W,\MF\}$, a target $f$, training inputs $A$ with $D=f|_A$, and a test input $x^*$: $e(x^*):=1-\nu(f(x^*)\mid D,x^*)$ is the posterior mass wrong or undefined at $x^*$ ($0$ if $x^*\in A$); $e^{\rm wr}(x^*):=\sum_{y\ne f(x^*)}\nu(D+(x^*,y))/\nu(D)$ is the mass defined and wrong; $e_{\rm norm}(x^*)$ is the error of the normalised predictor. Since the normaliser is at most $\nu(D)$, $e^{\rm wr}\le e_{\rm norm}\le e$: lower bounds on $e^{\rm wr}$ and upper bounds on $e$ hold in both conventions.

\begin{proposition}[random splits]\label{prop:learn:splits}
\status{proved; technique known; computed (context c3)}\src{context Prop 2.9}
Let $\nu\in\{W,\MF\}$, $B$ any finite input set (possibly adversarial), $f$ any labelling of $B$. Train on the first $k-1$ inputs of a uniformly random order of $B$ and test on the $k$-th, with $k$ uniform on $\{1,\dots,|B|\}$ and independent. The expected test loss is $-\log(\nu(f|_B)/\nu(\emptyset))/|B|\le L_\nu(f|_B)/|B|\le(\KF(f|_B)+c)/|B|$, with $c=0$ for $W$, and the expected test error is at most $(\ln2)L_\nu(f|_B)/|B|$.
\end{proposition}

\begin{proposition}[leave-one-out]\label{prop:learn:loo}
\status{proved; computed (context c7)}\src{context Prop 2.10}
For $\nu\in\{W,\MF\}$, a labelling $f$ of a finite $B$, and $e_{\rm LOO}(x):=1-\nu(f|_B)/\nu(f|_{B\setminus\{x\}})$: $\sum_{x\in B}e_{\rm LOO}(x)\le(\nu(\emptyset)-\nu(f|_B))/\nu(f|_B)<2^{L_\nu(f|_B)}\le2^{\KF(f|_B)+c}$.
\end{proposition}

\emph{Proof ideas} (\cref{app:learn:splits}). Splits: in every order the step losses sum to the order-free total. Leave-one-out: the hypotheses correct on $B\setminus\{x\}$ and not at $x$ form disjoint sets for different $x$.

\begin{conjecture}[leave-one-out for SFI]\label{conj:learn:loo}
\status{conjecture}\src{context Conj 2.11; context §7 item 1}
For SFI, $\frac1{|B|}\sum_{x\in B}e_{\rm LOO}(x)=O(\KF(f|_B)/|B|)$, up to logarithmic factors.
\end{conjecture}

For general prior weights \cref{prop:learn:loo} is tight within a factor of about~$2$ (\cref{rem:learn:increasing}; heuristic in \cref{app:learn:remarks}).

\begin{proposition}[separating properties are typical even in distribution]\label{prop:learn:typical}
\status{proved; computed (context c4)}\src{context Prop 2.12}
Let $f\equiv0$ on $\{0,1\}^d$, $x_1,\dots,x_n$ i.i.d.\ uniform training inputs and $x^*$ an independent uniform test input, $d\ge\ell:=\lceil\log n\rceil+1$. Let $G_u$ be the set of strings that start with the $\ell$-bit prefix $u$ of $x^*$. Then (1) $G_u$ separates $x^*$ from the training inputs with probability $(1-2^{-\ell})^n\ge\frac12$; (2) $K(G_u)\le\ell+2\log(\ell+1)+c$; (3) for $\nu\in\{W,\MF\}$, $\Ex\,e^{\rm wr}(x^*)\ge\frac12\cdot2^{-\ell-2\log(\ell+1)-c}\ge2^{-c}/(8n(\log n+3)^2)$, in both conventions. For computable binary $f$ the same argument gives $\Ex\,e^{\rm wr}(x^*)\ge2^{-K(f)-c}/(8n(\log n+3)^2)$.
\end{proposition}

\emph{Proof idea.} On the event of (1), answering~$1$ on $G_u$ and running a training-consistent program elsewhere costs $K(G_u)+c$ extra bits (a case of \cref{prop:ood:hanni}; \cref{app:learn:typical}). So ``no simple property separates test from train'' typically holds only at scales of about $\log n$ bits. Reading this as ``matching the Ces\`aro bound up to $\log^2n$'' is refuted (\cref{rem:learn:cesaro}). Per step:

\begin{corollary}[per-step in-distribution error]\label{cor:learn:perstep}
\status{proved; computed (context c7)}\src{context Cor 2.17}
Let $\mu$ be any distribution on $X$, $f$ partial computable with $\mu(\dom f)=1$, $x_1,\dots,x_n,x^*$ i.i.d.\ from $\mu$, $\nu\in\{W,\MF\}$, and $e$ the semimeasure-convention error (so the bounds hold normalised too).
(a) For every $\varepsilon\in(0,1)$, $\Ex e(x^*)\le\varepsilon+2^{K(f)+c}e^{-\varepsilon n}$; hence $\Ex e(x^*)\le((K(f)+c)\ln2+\ln n+1)/n$. With probability at least $1-\delta$ over the training sample, $\Ex_{x^*}[e(x^*)\mid D_n]\le\varepsilon+2^{K(f)+c}e^{-\varepsilon n}/\delta$. The posterior-majority predictor has expected 0--1 risk at most $2\Ex e(x^*)$.
(b) $\Ex e(x^*)\le\Ex\big[\frac{\nu(\emptyset)-\nu(f|_S)}{(n+1)\nu(f|_S)}\big]\le2^{K(f)+c}/(n+1)$, where $S$ is the set of the $n+1$ sampled values.
For $f\equiv0$ and uniform inputs on $\{0,1\}^d$, $d\ge\lceil\log n\rceil+1$: $2^{-c}/(8n(\log n+3)^2)\le\Ex e(x_{n+1})\le2^{c}/(n+1)$ in both conventions, so the per-step error is $\Theta(1/n)$ up to the $\log^2n$ factor.
\end{corollary}

\emph{Proof idea} (\cref{app:learn:typical}). (a) A hypothesis with error above $\varepsilon$ survives with probability at most $e^{-\varepsilon n}$; the normaliser is at least $2^{-K(f)-c}$. (b) By exchangeability the error is an average of leave-one-out errors. Part (a) is the Markov half of the slides' relevance-based bound \citep{vaintrobhanni2026nnbayes} for $W$ and $\MF$ (\cref{rem:learn:relations}).

% ------------------------------------------------------------------
\subsection{Comparison, and the sense of ``ideal''}\label{sec:learn:ideal}

\begin{proposition}[Occam bound in description-length form]\label{prop:learn:occam}
\status{proved (standard); references known, unverified}\src{foundations §6.6; F32}
Let $x_t$ be i.i.d.\ from $\mu$, $\varphi$ partial computable with $\mu(\dom\varphi)=1$, $y_t=\varphi(x_t)$. With probability at least $1-\delta$, every program $p\in\dom V$ consistent with $D_n$ has $\mu\{x:\varphi_{V(p)}(x)\ne\varphi(x)\}\le(|p|\ln2+\ln(1/\delta))/n$, where ``$\ne$'' includes ``undefined''. So the shortest consistent program, of length $\KF(D_n)\le K(\varphi)$, has 0--1 risk at most $(K(\varphi)\ln2+\ln(1/\delta))/n$.
\end{proposition}

So for deterministic targets SFI (\cref{thm:learn:batch}(a)) and the Kolmogorov version learn at the same rate, though their code lengths can differ (\cref{thm:versions:nocoding}). \Cref{app:learn:remarks} compares with Occam, PAC-Bayes, resolvability and expert-advice bounds \citep{blumer1987occam,mcallester1999some,barron1991minimum,vovk1990aggregating,cesabianchi2006prediction}.

\begin{remark}[in what sense SFI is ideal supervised learning]\label{rem:learn:ideal}
\status{remark (each item cited)}\src{foundations §6.7}
Established: (1) SFI is the Bayes predictor for the prior ``a random lower-semicomputable labeller'', machine laws weighted by $\mm(i)\ge2^{-K(i)}$ (\cref{thm:setting:universal}), and it is universal for that class pathwise (\cref{thm:learn:dominance}). (2) With i.i.d.\ inputs from any distribution and computable labels, the averaged predictor has expected log and 0--1 risk at most $(K(\varphi)+c)/n$; with frozen noise, or fresh noise on stamped inputs or for $\Mpq$, the cumulative bound $K(P)+c$ holds (\cref{thm:learn:batch,prop:learn:perquery}). This is the Occam rate with $K$ for $\log|H|$, without a union bound or confidence parameter. (3) Inputs are never predicted, so the input distribution may be arbitrary at no cost (chosen inputs: \cref{prop:learn:active}; positive examples and interventions: \cref{rem:learn:positive,rem:learn:interventions}).
Limits: \emph{loss}: the guarantees are for log loss on labels, and 0--1 bounds follow ($K/n$ without noise, $\sqrt{K/n}$ excess with noise); \emph{averaged predictor}: \cref{prop:learn:curve,rem:learn:lastiterate}; \emph{function semantics}: \cref{prop:learn:freshnoise}; \emph{computability}: \cref{prop:setting:computability}, and \cref{sec:comp} for computable versions; \emph{constants}: machine constants, and $K(\alpha,\beta)$ for encodings (\cref{thm:setting:invariance}); \emph{distribution shift}: \cref{sec:ood}; \emph{comparator class}: simple history-aware predictors have loss $0$ where SFI's loss is unbounded (\cref{thm:hist:chain,thm:hist:skipped,thm:ex:cycle}; the swap data of \cref{ex:setting:swap} do not show this), and \SFIhist{} removes this limit (\cref{cor:hist:H4}).
\end{remark}

\begin{remark}[prior work on the framing]\label{rem:learn:hutterEX}
\status{known; unverified}\src{foundations §6.6 last item; history §4.6 item 2}
A table-of-contents snippet shows a section ``Supervised Learning from Examples (EX)'' in \citet{hutter2007aixi}, and the abstracts of \citet{hutter2000theory} and of that chapter list supervised learning among the AIXI model's problem classes; \citet{hutter2005universal} is the book-length treatment. We have not read these sections; the universal mixture as ideal learner from examples is prior work in substance.
\end{remark}

